% Methodology section for EpiContext paper
\section{Methodology}

We present \EpiContext, a framework that reformulates agent context management through the lens of biological epigenetics. Figure~\ref{fig:architecture} provides an architectural overview. We first formalize the problem setting (Section~\ref{sec:problem}), then introduce the context graph representation (Section~\ref{sec:graph}), the three epigenetic operators (Section~\ref{sec:operators}), the fitness-driven evolutionary engine (Section~\ref{sec:fitness}), and finally the complete execution pipeline (Section~\ref{sec:pipeline}).

\begin{figure}[t]
\centering
% EpiContext Architecture Diagram (TikZ)
% Figure 1: System overview showing the epigenetic context regulation pipeline

\begin{tikzpicture}[
    node distance=0.8cm,
    box/.style={rectangle, draw=black!70, rounded corners=3pt, minimum width=2.2cm, minimum height=0.7cm, align=center, font=\footnotesize},
    bigbox/.style={rectangle, draw=black!40, rounded corners=5pt, minimum width=3.0cm, minimum height=2.8cm, fill=black!3},
    arrow/.style={->, >=stealth, thick, black!60},
    label/.style={font=\footnotesize\bfseries},
    smalllabel/.style={font=\scriptsize, text=black!60},
]

% ===== Left: Context Graph (Genome) =====
\node[bigbox] (genome) at (0, 0) {};
\node[label, above=0.1cm of genome.north] {Context Graph (Genome)};

\node[box, fill=blue!8] (sys) at (0, 0.8) {System\\Chromosome};
\node[box, fill=green!8] (env) at (0, 0.0) {Environment\\Chromosome};
\node[box, fill=orange!8] (tool) at (0, -0.8) {Tool\\Chromosome};
\node[box, fill=red!8] (mem) at (0, -1.6) {Memory\\Chromosome};

% ===== Middle: Epigenetic Operators =====
\node[bigbox, minimum width=3.5cm, minimum height=3.5cm] (ops) at (4.5, 0) {};
\node[label, above=0.1cm of ops.north] {Epigenetic Operators};

\node[box, fill=yellow!15, minimum width=2.8cm] (methyl) at (4.5, 0.9) {Methylation $\mathcal{M}$\\{\small Silence noisy blocks}};
\node[box, fill=cyan!12, minimum width=2.8cm] (acetyl) at (4.5, 0.0) {Acetylation $\mathcal{A}$\\{\small Activate relevant tools}};
\node[box, fill=purple!10, minimum width=2.8cm] (cross) at (4.5, -0.9) {Crossover $\mathcal{X}$\\{\small Explore variants}};

% ===== Right: Fitness + Payload =====
\node[bigbox, minimum width=3.2cm] (fitness) at (9.0, 0.5) {};
\node[label, above=0.1cm of fitness.north] {Fitness Function};

\node[font=\footnotesize, align=center] (feq) at (9.0, 0.5) {$\mathcal{F}(P) = \alpha\mathcal{R}_{\text{task}}$\\${}-\beta\mathcal{C}_{\text{token}}$\\${}+\gamma\mathcal{I}_{\text{density}}$};

\node[box, fill=green!8, minimum width=2.8cm] (payload) at (9.0, -1.5) {Optimized\\Payload $P_t^*$};

% ===== Arrows =====
\draw[arrow] (genome.east) -- (ops.west) node[midway, above, smalllabel] {$\mathcal{G}_t$};
\draw[arrow] (ops.east) -- (fitness.west) node[midway, above, smalllabel] {$\mathcal{V}_{\text{active}}$};
\draw[arrow] (fitness.south) -- (payload.north) node[midway, right, smalllabel] {$\arg\max\mathcal{F}$};

% ===== Feedback loop =====
\draw[arrow, dashed, black!40] (payload.south) -- ++(0, -0.6) -| (genome.south)
    node[pos=0.25, below, smalllabel] {weight update $\Delta w$};

% ===== Bottom: LLM =====
\node[box, fill=black!5, minimum width=3cm, minimum height=0.6cm] (llm) at (9.0, -2.8) {LLM Inference};
\draw[arrow] (payload.south) -- (llm.north);

% ===== Legend =====
\node[smalllabel, anchor=north west] at (-1.5, -3.5) {
    \textbf{Genome}: Complete interaction history (immutable)\\[2pt]
    \textbf{Epigenetic Expression}: Active context subset (dynamic)\\[2pt]
    \textbf{Fitness}: Joint optimization of task success \& token efficiency
};

\end{tikzpicture}
\caption{Architectural overview of \EpiContext. The complete interaction history is stored as an immutable Context Graph (genome) organized into four functional chromosomes. Three epigenetic operators---methylation (silencing), acetylation (activation), and crossover (exploration)---dynamically regulate which subset of the graph is expressed in each LLM request. A fitness function $\mathcal{F}(P) = \alpha\mathcal{R}_{\text{task}} - \beta\mathcal{C}_{\text{token}} + \gamma\mathcal{I}_{\text{density}}$ jointly optimizes for task success and token efficiency, with feedback updating epigenetic weights after each turn.}
\label{fig:architecture}
\end{figure}

\subsection{Problem Formulation}
\label{sec:problem}

Consider an LLM agent executing a task $\tau$ over $T$ turns. At each turn $t \in \{1, \ldots, T\}$, the agent produces a thought $h_t$, selects an action $a_t$ from a tool set $\mathcal{T} = \{f_1, \ldots, f_K\}$, and receives an observation $o_t$ from the environment. The agent's complete interaction history up to turn $t$ is:

\begin{equation}
    \mathcal{H}_t = \{(h_1, a_1, o_1), \ldots, (h_t, a_t, o_t)\}
\end{equation}

In standard agent architectures, the request payload $P_t$ sent to the LLM at turn $t+1$ is constructed as:

\begin{equation}
    P_t^{\text{standard}} = \mathcal{C}_{\text{sys}} \oplus \mathcal{C}_{\text{tool}} \oplus \mathcal{H}_t
\end{equation}

where $\mathcal{C}_{\text{sys}}$ is the system prompt, $\mathcal{C}_{\text{tool}}$ contains all $K$ tool schemas, and $\oplus$ denotes concatenation. As $t$ grows, $|P_t^{\text{standard}}|$ grows linearly with $t$, quickly exhausting context budgets and degrading model performance.

The core problem is to construct an optimized payload $P_t^*$ that maximizes task success while minimizing token consumption:

\begin{equation}
    P_t^* = \arg\max_{P \subseteq \mathcal{G}_t} \mathcal{F}(P)
\end{equation}

where $\mathcal{G}_t$ is the complete context graph (the ``genome'') and $\mathcal{F}$ is a fitness function defined in Section~\ref{sec:fitness}.

\subsection{Context Graph Representation}
\label{sec:graph}

\EpiContext{} replaces the flat message array with a structured \textbf{Context Graph} $\mathcal{G} = (\mathcal{V}, \mathcal{E})$, a directed acyclic graph (DAG) that serves as the agent's complete ``genome.''

\begin{definition}[Context Graph]
A context graph $\mathcal{G} = (\mathcal{V}, \mathcal{E})$ consists of:
\begin{itemize}
    \item \textbf{Nodes} $\mathcal{V}$: Each node $v_i \in \mathcal{V}$ represents a discrete event with attributes:
    \begin{itemize}
        \item Type $\tau(v_i) \in \{\text{system}, \text{thought}, \text{action}, \text{observation}, \text{error}, \text{summary}\}$
        \item Content $c(v_i)$: the textual content of the event
        \item Epigenetic tag $w(v_i) \in [0, 1]$: the current activation weight
        \item Timestamp $t(v_i)$: creation time
    \end{itemize}
    \item \textbf{Edges} $\mathcal{E}$: Directed edges $e_{ij} = (v_i, v_j)$ representing temporal or causal relationships between events.
\end{itemize}
\end{definition}

The epigenetic tag $w(v_i)$ is the central mechanism: $w(v_i) = 1$ indicates full activation (the node's content is included in the payload), while $w(v_i) = 0$ indicates complete silencing (methylation). Intermediate values represent partial activation.

\subsubsection{Chromosome Decomposition}

For efficient regulation, we decompose the context graph into four functional ``chromosomes,'' each representing a distinct category of information:

\begin{align}
    \mathcal{C}_{\text{sys}} &= \{v \in \mathcal{V} : \tau(v) = \text{system}\} \label{eq:csys} \\
    \mathcal{C}_{\text{env}} &= \{v \in \mathcal{V} : \tau(v) \in \{\text{file\_snapshot}, \text{terminal}\}\} \label{eq:cenv} \\
    \mathcal{C}_{\text{tool}} &= \{v \in \mathcal{V} : \tau(v) \in \{\text{tool\_schema}, \text{tool\_call}, \text{tool\_result}\}\} \label{eq:ctool} \\
    \mathcal{C}_{\text{mem}} &= \{v \in \mathcal{V} : \tau(v) \in \{\text{thought}, \text{action}, \text{observation}, \text{error}\}\} \label{eq:cmem}
\end{align}

This decomposition enables targeted regulation: tool acetylation operates primarily on $\mathcal{C}_{\text{tool}}$, while memory methylation targets $\mathcal{C}_{\text{mem}}$.

\subsection{Epigenetic Operators}
\label{sec:operators}

\EpiContext{} defines three epigenetic operators that dynamically regulate context graph expression. Each operator is designed to be computationally lightweight, requiring at most a single auxiliary LLM call for relevance assessment.

\subsubsection{Memory Methylation ($\mathcal{M}$)}

Methylation silences noisy or resolved interaction blocks. When the agent has resolved a subproblem after extensive trial-and-error (e.g., 10 turns to fix a dependency conflict), those detailed logs become noise for subsequent tasks.

\begin{definition}[Methylation Operator]
Given a set of node indices $I \subseteq \{1, \ldots, |\mathcal{V}|\}$ to silence, the methylation operator $\mathcal{M}$:
\begin{enumerate}
    \item Sets $w(v_i) \leftarrow 0$ for all $i \in I$ (or applies gradual decay: $w(v_i) \leftarrow w(v_i) \cdot \lambda$ with $\lambda \in (0, 1)$)
    \item Creates a summary node $v_s$ with $c(v_s) = \text{Summarize}(\{c(v_i)\}_{i \in I})$ and $w(v_s) = 1$
    \item Adds edges $(v_s, v_i)$ for all $i \in I$ with type ``summarizes''
\end{enumerate}
\end{definition}

Methylation is triggered when $|\{v \in \mathcal{V} : w(v) > \theta\}| > N_{\text{max}}$, where $\theta$ is the activation threshold and $N_{\text{max}}$ is the maximum desired active node count. The summarization can be performed by a lightweight LLM call or through extractive methods.

\subsubsection{Tool Acetylation ($\mathcal{A}$)}

Modern agents often carry 20+ tools, but any specific task typically requires only 2--3. Tool acetylation dynamically selects the relevant subset.

\begin{definition}[Acetylation Operator]
Given the tool chromosome $\mathcal{C}_{\text{tool}}$ and current task context $\tau$, the acetylation operator $\mathcal{A}$:
\begin{enumerate}
    \item Computes relevance scores $r(f_k, \tau)$ for each tool $f_k$ using multi-level matching:
    \begin{equation}
        r(f_k, \tau) = \lambda_1 \cdot \text{name\_match}(f_k, \tau) + \lambda_2 \cdot \text{desc\_match}(f_k, \tau) + \lambda_3 \cdot \text{param\_match}(f_k, \tau)
    \end{equation}
    \item Retains tools with $r(f_k, \tau) \geq \rho_{\text{min}}$, up to a maximum of $K_{\text{max}}$ tools
    \item Sets $w(v) \leftarrow 1$ for retained tool nodes and $w(v) \leftarrow 0$ for filtered ones
\end{enumerate}
\end{definition}

The relevance computation uses lexical overlap between tool names/descriptions/parameters and the task description. For more sophisticated matching, a lightweight embedding similarity or LLM-based relevance scoring can be employed.

\subsubsection{Contextual Crossover ($\mathcal{X}$)}

When the agent encounters deadlock---defined as $E$ consecutive failed actions---crossover generates and evaluates multiple context variants in parallel.

\begin{definition}[Crossover Operator]
Given the current active node set $\mathcal{V}_{\text{active}} = \{v \in \mathcal{V} : w(v) > \theta\}$, the crossover operator $\mathcal{X}$:
\begin{enumerate}
    \item Generates $M$ variant context sets $\{\mathcal{V}^{(1)}, \ldots, \mathcal{V}^{(M)}\}$ through stochastic mutation:
    \begin{itemize}
        \item \textbf{Dropout}: Randomly exclude nodes with probability $p_{\text{drop}}$
        \item \textbf{Resurrection}: Re-activate previously methylated nodes with probability $p_{\text{res}}$
        \item \textbf{Permutation}: Reorder memory nodes to test different temporal emphases
    \end{itemize}
    \item Evaluates each variant's fitness $\mathcal{F}(\mathcal{V}^{(m)})$
    \item Selects the best variant $\mathcal{V}^* = \arg\max_m \mathcal{F}(\mathcal{V}^{(m)})$
    \item Propagates: increases $w(v)$ by $\Delta w$ for nodes in $\mathcal{V}^*$
\end{enumerate}
\end{definition}

Crossover implements a form of parallel search over context configurations, analogous to how genetic recombination explores the fitness landscape in biological evolution.

\subsection{Fitness-Driven Evolution}
\label{sec:fitness}

The fitness function $\mathcal{F}$ provides the objective that guides all epigenetic regulation. It jointly optimizes three competing desiderata:

\begin{equation}
    \mathcal{F}(P) = \alpha \cdot \mathcal{R}_{\text{task}}(P) - \beta \cdot \mathcal{C}_{\text{token}}(P) + \gamma \cdot \mathcal{I}_{\text{density}}(P)
    \label{eq:fitness}
\end{equation}

where:

\begin{itemize}
    \item \textbf{Task Reward} $\mathcal{R}_{\text{task}}(P)$: Measures whether the payload led to successful task progress. Defined as:
    \begin{equation}
        \mathcal{R}_{\text{task}}(P) = \mathbf{1}[\text{task\_success}] + \eta \cdot \mathbf{1}[\text{tool\_call\_success}] - \epsilon \cdot n_{\text{errors}}
    \end{equation}
    where $\eta$ rewards successful tool calls and $\epsilon$ penalizes errors.
    
    \item \textbf{Token Cost} $\mathcal{C}_{\text{token}}(P)$: Normalized token count of the payload:
    \begin{equation}
        \mathcal{C}_{\text{token}}(P) = \frac{|P|}{B}
    \end{equation}
    where $B$ is a baseline token count (e.g., 10,000) for normalization.
    
    \item \textbf{Information Density} $\mathcal{I}_{\text{density}}(P)$: Ratio of effective decisions to total tokens:
    \begin{equation}
        \mathcal{I}_{\text{density}}(P) = \frac{n_{\text{effective}}}{n_{\text{total}}} \cdot \kappa
    \end{equation}
    where $n_{\text{effective}}$ counts actions that advanced the task, $n_{\text{total}}$ is the total token count, and $\kappa$ is a scaling factor.
\end{itemize}

The hyperparameters $(\alpha, \beta, \gamma)$ control the trade-off between task performance and efficiency. Higher $\alpha$ prioritizes task success; higher $\beta$ aggressively minimizes token usage; higher $\gamma$ favors information-dense contexts.

\subsubsection{Evolutionary Weight Update}

After each turn, epigenetic tags are updated based on the observed fitness:

\begin{equation}
    w(v_i)^{(t+1)} = \text{clip}\left(w(v_i)^{(t)} + \Delta w_i^{(t)}, 0, 1\right)
\end{equation}

where the update $\Delta w_i^{(t)}$ depends on the node's contribution to the turn outcome:

\begin{equation}
    \Delta w_i^{(t)} = \begin{cases}
        +\delta_{\text{reward}} & \text{if } v_i \text{ was active and turn succeeded} \\
        -\delta_{\text{penalize}} & \text{if } v_i \text{ was active and turn failed} \\
        +\delta_{\text{explore}} & \text{if } v_i \text{ was inactive (exploration bonus)}
    \end{cases}
\end{equation}

This update rule implements a simple reinforcement learning dynamic: nodes that participate in successful turns are reinforced, while those associated with failures are suppressed. The exploration bonus $\delta_{\text{explore}}$ prevents premature convergence by periodically reactivating silenced nodes.

\subsection{Execution Pipeline}
\label{sec:pipeline}

\EpiContext{} operates as a lightweight router between the agent client and the LLM gateway. Algorithm~\ref{alg:pipeline} describes the complete per-turn execution pipeline.

\begin{algorithm}[t]
\caption{\EpiContext{} Per-Turn Execution Pipeline}
\label{alg:pipeline}
\begin{algorithmic}[1]
\REQUIRE Context graph $\mathcal{G}$, task $\tau$, tool registry $\mathcal{T}$, fitness params $(\alpha, \beta, \gamma)$
\ENSURE Optimized payload $P_t$, updated graph $\mathcal{G}'$
\STATE \textbf{// Step 1: State Capture}
\STATE $v_h \leftarrow \text{AddNode}(\mathcal{G}, \text{``thought''}, h_t)$
\STATE $v_a \leftarrow \text{AddNode}(\mathcal{G}, \text{``action''}, a_t)$
\STATE $v_o \leftarrow \text{AddNode}(\mathcal{G}, \text{type}(o_t), o_t)$
\STATE $\text{AddEdge}(v_h, v_a)$, $\text{AddEdge}(v_a, v_o)$
\STATE \textbf{// Step 2: Epigenetic Masking}
\STATE $\mathcal{T}_{\text{active}} \leftarrow \mathcal{A}(\mathcal{C}_{\text{tool}}, \tau)$ \COMMENT{Acetylation}
\IF{$|\{v \in \mathcal{V} : w(v) > \theta\}| > N_{\text{max}}$}
    \STATE $I_{\text{old}} \leftarrow \text{SelectOldest}(\mathcal{V}_{\text{active}}, N_{\text{max}})$
    \STATE $\mathcal{M}(I_{\text{old}})$ \COMMENT{Methylation}
\ENDIF
\IF{$n_{\text{consecutive\_errors}} \geq E_{\text{threshold}}$}
    \STATE $\{\mathcal{V}^{(m)}\} \leftarrow \mathcal{X}(\mathcal{V}_{\text{active}}, M)$ \COMMENT{Crossover}
\ENDIF
\STATE \textbf{// Step 3: Payload Assembly}
\STATE $P_t \leftarrow \text{Assemble}(\mathcal{V}_{\text{active}}, \mathcal{T}_{\text{active}}, \tau)$
\STATE \textbf{// Step 4: LLM Inference}
\STATE $\text{response} \leftarrow \text{LLM}(P_t)$ \COMMENT{Send to LLM}
\STATE \textbf{// Step 5: Fitness Evaluation \& Update}
\STATE $f \leftarrow \mathcal{F}(P_t)$
\FORALL{$v \in \mathcal{V}_{\text{active}}$}
    \STATE $w(v) \leftarrow \text{clip}(w(v) + \Delta w(f, v), 0, 1)$
\ENDFOR
\RETURN $P_t, \mathcal{G}'$
\end{algorithmic}
\end{algorithm}

\subsubsection{Computational Overhead}

The primary computational cost of \EpiContext{} comes from the auxiliary LLM calls for summarization (methylation) and relevance assessment (acetylation). However, these calls use significantly fewer tokens than the savings they enable. Methylation summarization processes $N_{\text{methyl}}$ nodes but produces a single compact summary, with cost $O(N_{\text{methyl}} \cdot L_{\text{avg}})$ input tokens and $O(L_{\text{summary}})$ output tokens, where $L_{\text{avg}}$ is average node content length and $L_{\text{summary}} \ll N_{\text{methyl}} \cdot L_{\text{avg}}$. Acetylation relevance requires processing $K$ tool descriptions against the task, costing $O(K \cdot L_{\text{tool}} + L_{\text{task}})$ tokens, which is typically negligible compared to per-turn context costs. Crossover is only triggered on deadlock (fewer than 5\% of turns in practice); each variant evaluation costs one LLM call, but the parallel exploration often resolves deadlocks faster than sequential retry. In our experiments, the auxiliary LLM overhead accounts for less than 3\% of total token consumption while enabling 30--70\% reduction in primary context tokens.

\subsubsection{Theoretical Properties}

\begin{theorem}[Convergence of Epigenetic Weights]
Under the update rule in Equation~(10), with $\delta_{\text{reward}} > 0$ and $\delta_{\text{penalize}} > 0$, the epigenetic weights $w(v_i)$ converge to a stationary distribution where nodes that consistently contribute to successful turns have $w \to 1$ and nodes associated with failures have $w \to 0$, provided the task distribution is stationary.
\end{theorem}

\begin{proof}[Sketch]
The weight update is a bounded random walk with drift. For a node $v_i$, let $p_i$ be the probability that activating $v_i$ leads to a successful turn. The expected drift is $\mathbb{E}[\Delta w_i] = p_i \delta_{\text{reward}} - (1-p_i) \delta_{\text{penalize}}$. When $p_i > \frac{\delta_{\text{penalize}}}{\delta_{\text{reward}} + \delta_{\text{penalize}}}$, the drift is positive and $w_i \to 1$; otherwise $w_i \to 0$. The clipping to $[0,1]$ ensures boundedness. Full proof in Appendix~\ref{sec:proof}.
\end{proof}

This theorem establishes that \EpiContext{} naturally learns to activate useful context elements and suppress harmful ones, without requiring explicit supervision or gradient-based optimization.